% Luc Destombe --- licence CC BY-NC-SA 4.0
% https://creativecommons.org/licenses/by-nc-sa/4.0/deed.fr
% Fichier autonome : compiler avec pdflatex, deux fois (signets, nombre de pages).
\def\lyceeversion{1.0}
\documentclass[11pt,a4paper]{article}
\makeatletter
\newif\ifeleve
\elevefalse

\newif\iflycee@cm
\newif\iflycee@palatino

\RequirePackage[utf8]{inputenc}
\RequirePackage[T1]{fontenc}
\RequirePackage[french]{babel}
\RequirePackage{etoolbox}

\newcommand{\ShorthandsOff}{\@ifundefined{shorthandoff}{}{\shorthandoff{;:!?}}}
\newcommand{\ShorthandsOn}{\@ifundefined{shorthandon}{}{\shorthandon{;:!?}}}
\ShorthandsOff
\AtBeginDocument{\ShorthandsOn}
\AtBeginEnvironment{tikzpicture}{\ShorthandsOff}

\RequirePackage{geometry}
\RequirePackage{fancyhdr}
\RequirePackage{lastpage}
\RequirePackage{multicol}
\RequirePackage{array}
\RequirePackage{enumitem}
\RequirePackage{xcolor}
\RequirePackage{needspace}

\RequirePackage{amsmath,amssymb}
\iflycee@cm\else\RequirePackage{fourier}\fi
\DeclareMathAlphabet{\mathcal}{OMS}{cmsy}{m}{n}
\renewcommand{\ttdefault}{lmtt}
\begingroup\catcode`\_=\active \catcode`\^=\active
\global\let_\sb \global\let^\sp
\endgroup
\AtBeginDocument{\catcode`\_=12 \mathcode`\_="8000
  \catcode`\^=12 \mathcode`\^="8000 }
\newlength{\retraitcalculs}
\setlength{\retraitcalculs}{2em}

\RequirePackage{tikz}
\usetikzlibrary{arrows.meta,calc,babel,shapes.misc}
\RequirePackage[most]{tcolorbox}
\tcbuselibrary{skins,breakable}

\definecolor{coulDef}{HTML}{1F4E79}
\definecolor{coulProp}{HTML}{7030A0}
\definecolor{coulRem}{HTML}{C55A11}
\definecolor{coulEx}{HTML}{2E7D32}
\definecolor{coulExo}{HTML}{B03A2E}
\definecolor{coulPart}{HTML}{1F4E79}
\definecolor{coulMeth}{HTML}{1B6E4F}
\definecolor{coulCourbe}{HTML}{0B5ED7}
\definecolor{coulCourbeB}{HTML}{C00000}
\definecolor{coulCourbeC}{HTML}{2E7D32}
\definecolor{coulGrille}{HTML}{8C8C8C}
\definecolor{coulRep}{HTML}{1B6E4F}
\definecolor{coulBar}{HTML}{2E7D32}

\newcommand{\R}{\mathbb{R}}
\newcommand{\Rs}{\mathbb{R}^{*}}
\newcommand{\N}{\mathbb{N}}
\newcommand{\Ns}{\mathbb{N}^{*}}
\newcommand{\Z}{\mathbb{Z}}
\newcommand{\Q}{\mathbb{Q}}
\renewcommand{\le}{\leqslant}
\renewcommand{\ge}{\geqslant}
\renewcommand{\approx}{\simeq}
\newcommand{\vect}[1]{\overrightarrow{#1}}
\newcommand{\norme}[1]{\left\lVert #1 \right\rVert}
\newcommand{\intervalle}[2]{\left[#1\,;#2\right]}
\RequirePackage{cancel}
\renewcommand{\CancelColor}{\color{coulExo}}
\newcommand{\fois}[1]{\times{\color{coulExo}#1}}
\newcommand{\barre}[1]{\cancel{#1}}
\newlength{\lycee@eqh}\newlength{\lycee@eqd}
\newsavebox{\lycee@eqa}\newsavebox{\lycee@eqb}\newsavebox{\lycee@eqc}
\newcommand{\lycee@eqcol}[1]{\hbox{$\renewcommand{\arraystretch}{1.35}%
  \begin{array}[t]{@{}>{\displaystyle}r@{}>{\displaystyle{}}l@{}}#1\end{array}$}}
\newcommand{\lycee@eqmesure}[1]{%
  \ifdim\ht#1>\lycee@eqh\setlength{\lycee@eqh}{\ht#1}\fi
  \ifdim\dp#1>\lycee@eqd\setlength{\lycee@eqd}{\dp#1}\fi}
\newcommand{\lycee@eqtrait}{\hspace{1.2em}ou\hspace{1.2em}}
\newcommand{\lycee@eqiff}{$\iff$\hspace{0.8em}}
\newcommand{\eqproduit}[2]{\par\smallskip\noindent C'est une équation produit :
  $A\times B=0 \iff A=0$ ou $B=0$.\par\nopagebreak\smallskip
  \begingroup\setlength{\lycee@eqh}{0pt}\setlength{\lycee@eqd}{0pt}%
  \setbox\lycee@eqa=\lycee@eqcol{#1}\setbox\lycee@eqb=\lycee@eqcol{#2}%
  \lycee@eqmesure\lycee@eqa\lycee@eqmesure\lycee@eqb%
  \noindent\hspace*{\retraitcalculs}\lycee@eqiff\usebox\lycee@eqa\lycee@eqtrait\usebox\lycee@eqb\endgroup\par\smallskip}
\newcommand{\eqproduittrois}[3]{\par\smallskip\noindent C'est une équation produit :
  $A\times B\times C=0 \iff A=0$ ou $B=0$ ou $C=0$.\par\nopagebreak\smallskip
  \begingroup\setlength{\lycee@eqh}{0pt}\setlength{\lycee@eqd}{0pt}%
  \setbox\lycee@eqa=\lycee@eqcol{#1}\setbox\lycee@eqb=\lycee@eqcol{#2}\setbox\lycee@eqc=\lycee@eqcol{#3}%
  \lycee@eqmesure\lycee@eqa\lycee@eqmesure\lycee@eqb\lycee@eqmesure\lycee@eqc%
  \noindent\hspace*{\retraitcalculs}\lycee@eqiff\usebox\lycee@eqa\lycee@eqtrait\usebox\lycee@eqb\lycee@eqtrait\usebox\lycee@eqc\endgroup
  \par\smallskip}

\tikzset{grille/.style={coulGrille,line width=0.4pt}}
\newcommand{\quadrillage}[4]{\draw[grille,step=1] (#1,#2) grid (#3,#4);}
\tikzset{
  croix/.style={cross out,draw,line width=0.6pt,inner sep=0pt,minimum size=4pt},
  croixdroite/.style={croix,rotate=45,line width=0.8pt},
}
\newcommand{\pt}[3]{\path (#1) node[croix]{} node[#2,font=\small,inner sep=2.5pt]{$#3$};}

\newcommand{\lycee@repere}[4]{%
  \draw[grille,step=1] (#1,#2) grid (#3,#4);
  \draw[-{Stealth[length=2mm]},thick] (#1,0)--({#3+0.4},0) node[below left=-1pt]{$x$};
  \draw[-{Stealth[length=2mm]},thick] (0,#2)--(0,{#4+0.4}) node[below left=-1pt]{$y$};
  \node[below left,font=\scriptsize,inner sep=1.5pt] at (0,0) {$0$};}
\newcommand{\lycee@gradx}[1]{\foreach \k/\l in {#1}{%
  \draw (\k,0.07)--(\k,-0.07) node[below,font=\scriptsize,inner sep=1.5pt]{$\l$};}}
\newcommand{\lycee@grady}[1]{\foreach \k/\l in {#1}{%
  \draw (0.07,\k)--(-0.07,\k) node[left,font=\scriptsize,inner sep=1.5pt]{$\l$};}}
\newcommand{\lycee@intff}[2]{\left[#1\,;#2\right]}
\newcommand{\lycee@intfo}[2]{\left[#1\,;#2\right[}
\newcommand{\lycee@intof}[2]{\left]#1\,;#2\right]}
\newcommand{\lycee@intoo}[2]{\left]#1\,;#2\right[}
\newcommand{\lycee@ens}[1]{\left\{#1\right\}}
\AtBeginDocument{%
  \providecommand{\repere}{\lycee@repere}%
  \providecommand{\gradx}{\lycee@gradx}%
  \providecommand{\grady}{\lycee@grady}%
  \providecommand{\Cf}{\mathcal{C}\sb{f}}%
  \providecommand{\Cg}{\mathcal{C}\sb{g}}%
  \providecommand{\intff}{\lycee@intff}%
  \providecommand{\intfo}{\lycee@intfo}%
  \providecommand{\intof}{\lycee@intof}%
  \providecommand{\intoo}{\lycee@intoo}%
  \providecommand{\ens}{\lycee@ens}}

\newlist{questions}{enumerate}{3}
\setlist[questions,1]{label=\arabic*\textdegree),leftmargin=*,itemsep=2pt,topsep=3pt}
\setlist[questions,2]{label=\alph*),leftmargin=*,itemsep=1pt,topsep=2pt}
\setlist[questions,3]{label=\roman*),leftmargin=*,itemsep=1pt,topsep=1pt}
\setlist[itemize]{label=\textbullet,leftmargin=*,itemsep=2pt,topsep=3pt}

\newcommand{\besoinplace}[1]{\par\penalty\@M\begingroup
  \dimen@=#1\relax\dimen@ii=\pagegoal\advance\dimen@ii-\pagetotal
  \ifdim\dimen@>\dimen@ii\ifdim\dimen@ii>\z@\vfil\fi\break\fi\endgroup}

\newcommand{\lycee@pied}{\thepage/\pageref{LastPage}}
\newcommand{\entete}[2]{%
  \pagestyle{fancy}\fancyhf{}%
  \fancyhead[L]{\footnotesize\color{coulPart}\textbf{#1}}%
  \fancyhead[R]{\footnotesize\color{coulPart}\textbf{#2}}%
  \fancyfoot[C]{\footnotesize\lycee@pied}%
  \renewcommand{\headrulewidth}{0.4pt}%
  \renewcommand{\headrule}{\hbox to\headwidth{%
    \color{coulPart!40}\leaders\hrule height \headrulewidth\hfill}}}

\RequirePackage{ccicons}
\newcommand{\lycee@licence}{{\scriptsize\color{gray}%
  \copyright\ Luc Destombe --- \ccbyncsa\ CC BY-NC-SA 4.0}}
\newif\iflycee@licence \lycee@licencetrue
\newcommand{\sanslicence}{\lycee@licencefalse}
\AtBeginDocument{\iflycee@licence\AddToHookNext{shipout/foreground}{%
  \put(0,0){\rlap{\kern\dimexpr1in+\hoffset+\oddsidemargin+\textwidth\relax
    \llap{\raisebox{-\dimexpr1in+\voffset+\topmargin+\headheight+\headsep
      +\textheight+\footskip\relax}{\lycee@licence}}}}}\fi}

\newcommand{\lycee@site}{faire-des-maths.fr}
\newcommand{\lycee@version}{{\scriptsize\color{gray}\lycee@site
  \ifdefined\lyceeversion\ --- Version~\lyceeversion\fi}}
\AtBeginDocument{\iflycee@licence\AddToHookNext{shipout/foreground}{%
  \put(0,0){\rlap{\kern\dimexpr1in+\hoffset+\oddsidemargin\relax
    \raisebox{-\dimexpr1in+\voffset+\topmargin+\headheight+\headsep
      +\textheight+\footskip\relax}{\lycee@version}}}}\fi}

\newcommand{\formulesserrees}{%
  \setlength{\abovedisplayskip}{3pt}\setlength{\belowdisplayskip}{3pt}%
  \setlength{\abovedisplayshortskip}{2pt}\setlength{\belowdisplayshortskip}{3pt}}

\setlength{\parindent}{0pt}

\RequirePackage[implicit=false,hidelinks,pdfencoding=auto,psdextra,bookmarksopen,
  bookmarksopenlevel=1,pdfcreator={},pdfproducer={}]{hyperref}
\RequirePackage{bookmark}
\hypersetup{pdfauthor={Luc Destombe},pdfkeywords={CC BY-NC-SA 4.0}}
\newcounter{lycee@signet}
\newcommand{\signet}[3][]{\stepcounter{lycee@signet}%
  \edef\lycee@dest{\if\relax\detokenize{#1}\relax
    signet.\arabic{lycee@signet}\else#1\fi}%
  \hypertarget{\lycee@dest}{}\bookmark[level=#2,dest=\lycee@dest]{#3}}
\pdfstringdefDisableCommands{%
  \def\R{R}\def\N{N}\def\Z{Z}\def\Q{Q}\def\vect#1{#1}\def\Cf{Cf}\def\Cg{Cg}\def\fois#1{×#1}%
  \def\barre#1{#1}\def\intervalle#1#2{[#1;#2]}\def\intff#1#2{[#1;#2]}%
  \def\textsuperscript#1{#1}\def\dfrac#1#2{#1/#2}\def\frac#1#2{#1/#2}%
  \def\vec#1{#1⃗}}%
\newcommand{\res}[1]{#1}
\geometry{top=2cm,bottom=1cm,left=1cm,right=1cm,headheight=14pt,footskip=0.6cm}

\RequirePackage{pgfplots}
\pgfplotsset{compat=1.18}
\RequirePackage{tkz-tab}

\newcounter{partiecours}
\renewcommand{\thepartiecours}{\Roman{partiecours}}
\newcommand{\partiecours}[1]{%
  \refstepcounter{partiecours}%
  \Needspace*{8\baselineskip}%
  \bigskip
  \noindent\begin{tcolorbox}[enhanced,colback=coulPart,colframe=coulPart,
      boxrule=0pt,arc=2pt,left=8pt,right=8pt,top=4pt,bottom=4pt,
      before skip=6pt,after skip=10pt]
    \color{white}\bfseries\large \thepartiecours{} --- #1%
    \signet[partie-\arabic{partiecours}]{1}{\thepartiecours{} --- #1}
  \end{tcolorbox}%
  \addcontentsline{toc}{section}{\thepartiecours{} --- #1}%
}

\newtcolorbox{boitecours}[3][]{%
  enhanced,breakable,
  colback=white,colframe=#2,coltitle=white,
  fonttitle=\bfseries,
  attach boxed title to top left={xshift=8pt,yshift=-\tcboxedtitleheight/2},
  boxed title style={colback=#2,arc=2pt,boxrule=0pt},
  boxrule=0.9pt,arc=2pt,
  left=8pt,right=8pt,top=8pt,bottom=6pt,
  title={#3},#1}

\newcounter{methode}
\newenvironment{definitionbox}[1][Définition]%
  {\begin{boitecours}[phantom={\signet{2}{#1}}]{coulDef}{#1}}{\end{boitecours}}
\newenvironment{proprietebox}[1][Propriété]%
  {\begin{boitecours}[phantom={\signet{2}{#1}}]{coulProp}{#1}}{\end{boitecours}}
\newenvironment{methodebox}[1][Méthode]{\stepcounter{methode}%
  \begin{boitecours}[phantom={\signet[methode-\arabic{methode}]{2}{#1}}]{coulMeth}{#1}}%
  {\end{boitecours}}

\newcommand{\etape}[1]{\par\smallskip\textbf{\color{coulMeth}#1}\par\nobreak\smallskip}
\newcommand{\enonce}[1]{\emph{#1}\par\smallskip}
\newcommand{\reponse}[1]{\par\smallskip
  {\footnotesize\itshape\color{black!55}Réponses : #1}}
\newcommand{\demo}[1]{\textbf{\color{coulMeth}Démonstration.} #1}
\newcommand{\afaire}[1]{%
  \par\smallskip
  \noindent{\small\color{coulExo}$\blacktriangleright$\ \textbf{À faire :}
    fiche d'exercices du chapitre, #1.}\par\smallskip}

\newenvironment{calculs}{\setlength{\jot}{5pt}\@fleqntrue\@mathmargin\retraitcalculs
  \setlength{\abovedisplayskip}{4pt}\setlength{\belowdisplayskip}{4pt}%
  \setlength{\abovedisplayshortskip}{4pt}\setlength{\belowdisplayshortskip}{4pt}%
  \csname alignat*\endcsname{2}}%
  {\csname endalignat*\endcsname}
\newcommand{\rem}[1]{\text{\small\itshape\color{black!60}#1}}
\newcommand{\explique}[2]{\par\smallskip\noindent
  \begin{minipage}[t]{0.57\linewidth}\raggedright #1\end{minipage}\hfill
  \begin{minipage}[t]{0.40\linewidth}\raggedright\leavevmode
    {\small\itshape\color{black!60}#2}\end{minipage}\par}
\newcommand{\cas}[1]{\par\smallskip\noindent\textbf{\color{coulExo}#1}\vadjust{\nobreak}}
\newcommand{\calc}[1]{\par\smallskip\textbf{\color{coulExo}#1}\par\nobreak}

\newtcolorbox{remarquebox}[1][Remarque]{%
  enhanced,breakable,colback=white,colframe=coulRem,
  boxrule=0pt,leftrule=2.5pt,arc=0pt,outer arc=0pt,
  left=8pt,right=6pt,top=5pt,bottom=5pt,
  before upper={\textbf{\color{coulRem}#1}\par\smallskip}}

\newtcolorbox{exemplebox}[1][Exemple]{%
  enhanced,breakable,colback=white,colframe=coulEx,
  boxrule=0pt,leftrule=2.5pt,arc=0pt,outer arc=0pt,
  left=8pt,right=6pt,top=5pt,bottom=5pt,
  before upper={\textbf{\color{coulEx}#1}\par\smallskip}}

\pgfplotsset{
  repere/.style={
    axis lines=middle,
    axis line style={-{Stealth[length=2.4mm]},thick},
    every axis x label/.style={at={(ticklabel* cs:1.0)},anchor=west},
    every axis y label/.style={at={(ticklabel* cs:1.0)},anchor=south},
    label style={font=\footnotesize},
    tick label style={font=\scriptsize},
    grid=both,
    major grid style={grille},
    minor tick num=0,
    tick align=inside,
    clip mode=individual,
  },
}
\tikzset{
  courbe/.style={coulCourbe,line width=1.1pt,smooth},
  courbeB/.style={coulCourbeB,line width=1.1pt,smooth},
  courbeC/.style={coulCourbeC,line width=1.1pt,smooth},
}

\newlist{etapes}{enumerate}{1}
\setlist[etapes]{label=\textbf{\arabic*.},leftmargin=*,itemsep=1pt,topsep=2pt}

\newlist{expressions}{enumerate}{1}
\setlist[expressions]{label={\Alph*\ $=$},leftmargin=*,itemsep=3pt,topsep=3pt}

\newcounter{exo}
\renewcommand{\theexo}{\ifnum\value{exo}<10 0\fi\arabic{exo}}
\newtcolorbox{exercice}[1][]{%
  enhanced,breakable,
  colback=white,colframe=coulExo,
  boxrule=0.9pt,arc=2pt,
  left=8pt,right=8pt,top=8pt,bottom=6pt,
  coltitle=white,fonttitle=\bfseries,
  attach boxed title to top left={xshift=8pt,yshift=-\tcboxedtitleheight/2},
  boxed title style={colback=coulExo,arc=2pt,boxrule=0pt},
  phantom={\signet{2}{Exercice \arabic{exo}}},
  title={Exercice \theexo\ifx\relax#1\relax\else{} --- #1\fi}}
\newenvironment{exo}[1][\relax]{\refstepcounter{exo}\begin{exercice}[#1]}{\end{exercice}}

  \newenvironment{correction}{\par}{\par}
  \newcommand{\autableau}{}
  \newcommand{\etapeexemples}{\etape{Exemples corrigés}}
  \newcommand{\etapetableau}[1]{\etape{#1}}
  \newenvironment{exempletableau}[1][Exemple]%
    {\begin{exemplebox}[#1]}{\end{exemplebox}}

\RequirePackage{comment}
\excludecomment{correctionavj}

\newcommand{\coursEntete}{}
\newcommand{\coursNiveau}{}
\newcommand{\coursTitre}{}
\newcommand{\coursSurTitre}{}
\newcommand{\coursSousTitre}{}

\newcommand{\enteteCours}[1]{\renewcommand{\coursEntete}{#1}}
\newcommand{\chapitrecours}[2]{\enteteCours{Chapitre #1 --- #2}}
\newcommand{\niveaucours}[1]{\renewcommand{\coursNiveau}{#1}}
\newcommand{\titrecours}[3][]{%
  \renewcommand{\coursTitre}{#2}%
  \renewcommand{\coursSurTitre}{#1}%
  \renewcommand{\coursSousTitre}{#3}}

\pagestyle{fancy}
\fancyhf{}
\fancyhead[L]{\footnotesize\color{coulPart}\textbf{\coursEntete}}
\fancyhead[R]{\footnotesize\color{coulPart}\textbf{\coursNiveau}}
\fancyfoot[C]{\footnotesize\thepage}
\renewcommand{\headrulewidth}{0.4pt}
\renewcommand{\headrule}{\hbox to\headwidth{\color{coulPart!40}\leaders\hrule height \headrulewidth\hfill}}

\setlength{\parskip}{2pt}

\newcommand{\lycee@mentioneleve}{\\[2pt]{\small\itshape Version élève}}
\AtBeginDocument{%
  \begin{center}
    \begin{tcolorbox}[enhanced,width=0.72\linewidth,colback=white,colframe=coulPart,
        boxrule=1.6pt,arc=3pt,halign=center,top=10pt,bottom=10pt]
      \ifx\coursSurTitre\empty
        {\Huge\bfseries\color{coulPart} \coursTitre}\\[4pt]
      \else
        {\Huge\bfseries\color{coulPart} \coursTitre}\\[3pt]
        {\Large\bfseries\color{coulPart} \coursSurTitre}\\[5pt]
      \fi
      {\normalsize \coursSousTitre}%
      
    \end{tcolorbox}
  \end{center}%
}

\RequirePackage{tikz}
\RequirePackage{xcolor}

\usetikzlibrary{decorations.text}
\definecolor{lycee@metal}{HTML}{1F4E79}
\definecolor{lycee@manche}{HTML}{8A5A2B}
\definecolor{lycee@etincelle}{HTML}{F28C28}
\newcommand{\lycee@pluriel}[2]{#1 #2\ifnum#1>1 s\fi}
\newcommand{\lycee@sceau}{%
  \edef\lycee@bas{\lycee@pluriel{\lyceerelectures}{relecture} et
    \lycee@pluriel{\lyceerodages}{rodage}}%
  \begin{tikzpicture}[line join=round]
    \fill[coulPart!10,draw=coulPart,line width=.8pt] (0:2.2)
      \foreach \i in {1,...,71}{ -- (\i*5:{isodd(\i)?2.1:2.2}) } -- cycle;
    \draw[coulPart,line width=1.2pt] (0,0) circle (1.93);
    \draw[coulPart,line width=.4pt] (0,0) circle (1.85);
    \draw[coulPart,line width=.4pt] (0,0) circle (1.2);
    \path[decorate,decoration={text along path,text align=center,
        text={|\fontsize{9.5}{11}\selectfont\bfseries\color{coulPart}|TEST{\'E} AVEC MES {\'E}L{\`E}VES}}]
      (180:1.45) arc (180:0:1.45);
    \path[decorate,decoration={text along path,text align=center,
        text/.expanded={|\noexpand\fontsize{8.5}{10}\noexpand\selectfont\noexpand\color{coulPart}|\lycee@bas}}]
      (195:1.68) arc (195:345:1.68);
    \begin{scope}[x=.25mm,y=-.25mm]
      \begin{scope}[shift={(-32,-40)}]%
        \fill[lycee@metal] (23,31) rectangle (41,49);
        \draw[lycee@etincelle,line width=.5pt,line cap=round]
          (20,30)--(15,27) (19,34)--(14,35) (44,30)--(49,27) (45,34)--(50,35);
        \fill[lycee@manche,rounded corners=.3mm] (38,24.6) rectangle (60,27.6);
        \fill[lycee@metal,rounded corners=.4mm] (25,20) rectangle (39,31);
        \fill[lycee@etincelle,rounded corners=.2mm] (25,20) rectangle (28,31);
      \end{scope}
    \end{scope}
    \ifdefined\lyceesceauplusieurs\def\lycee@annee{Ann{\'e}es~:}\else
      \def\lycee@annee{Ann{\'e}e}\fi
    \path[decorate,decoration={text along path,text align=center,
        text/.expanded={|\noexpand\fontsize{8}{10}\noexpand\selectfont\noexpand\color{coulPart}|%
          \unexpanded\expandafter{\lycee@annee} \unexpanded\expandafter{\lyceesceauannees}}}]
      (95:1.08) arc (95:445:1.08);
  \end{tikzpicture}}
\AtEndDocument{\ifdefined\lyceesceauannees
  \par\besoinplace{5cm}\vspace{1.5em}%
  \begin{center}\lycee@sceau\end{center}\fi}
\makeatother

\chapitrecours{7}{Loi binomiale}
\niveaucours{T\textsuperscript{le} spécialité}
\titrecours{LOI BINOMIALE}{Cours --- Terminale spécialité}

\begin{document}

\newtcblisting{python}{%
  listing only,enhanced,breakable,
  colback=black!3,colframe=black!35,boxrule=0.5pt,arc=1pt,
  left=6pt,right=6pt,top=2pt,bottom=2pt,
  listing options={language=Python,basicstyle=\ttfamily\small,
    keywordstyle=\color{coulMeth}\bfseries,columns=fullflexible,
    keepspaces=true,showstringspaces=false,aboveskip=0pt,belowskip=0pt}}

\tikzset{evt/.style={inner sep=2pt},
  proba/.style={font=\footnotesize,fill=white,inner sep=1pt},
  nom/.style={font=\footnotesize,inner sep=1.5pt,sloped}}
\newcommand{\arbfeuilles}[4]{\def\arbfAB{#1}\def\arbfAnB{#2}\def\arbfnAB{#3}\def\arbfnAnB{#4}}
\newcommand{\arbsansfeuilles}{\def\arbpasfeuilles{}}
\newcommand{\arbvaleurs}{\def\arbseules{}\arbsansfeuilles}
\newcommand{\arblab}[2]{\ifx\arbseules\undefined
  $#1\if\relax\detokenize{#2}\relax\else=#2\fi$\else$#2$\fi}
\newcommand{\arbfeuille}[2]{%
  \node[right,font=\footnotesize] at (#1.east) {$#2\ifx\arbtmp\empty\else=\arbtmp\fi$};}
\newcommand{\arbre}[8]{%
  \def\arbNA{\overline{#1}}\def\arbNB{\overline{#2}}%
  \node (O) at (0,0) {};
  \node[evt] (A) at (3.4,1.2) {$#1$};
  \node[evt] (nA) at (3.4,-1.2) {$\arbNA$};
  \node[evt] (AB) at (7.2,1.9) {$#2$};
  \node[evt] (AnB) at (7.2,0.5) {$\arbNB$};
  \node[evt] (nAB) at (7.2,-0.5) {$#2$};
  \node[evt] (nAnB) at (7.2,-1.9) {$\arbNB$};
  \draw (O.center) -- node[nom,above]{\arblab{P(#1)}{#3}} (A);
  \draw (O.center) -- node[nom,below]{\arblab{P(\arbNA)}{#4}} (nA);
  \draw (A) -- node[nom,above]{\arblab{P_{#1}(#2)}{#5}} (AB);
  \draw (A) -- node[nom,below]{\arblab{P_{#1}(\arbNB)}{#6}} (AnB);
  \draw (nA) -- node[nom,above]{\arblab{P_{\arbNA}(#2)}{#7}} (nAB);
  \draw (nA) -- node[nom,below]{\arblab{P_{\arbNA}(\arbNB)}{#8}} (nAnB);
  \ifx\arbpasfeuilles\undefined
    \ifx\arbfAB\undefined\arbfeuilles{}{}{}{}\fi
    \let\arbtmp\arbfAB\arbfeuille{AB}{P(#1\cap #2)}
    \let\arbtmp\arbfAnB\arbfeuille{AnB}{P(#1\cap\arbNB)}
    \let\arbtmp\arbfnAB\arbfeuille{nAB}{P(\arbNA\cap #2)}
    \let\arbtmp\arbfnAnB\arbfeuille{nAnB}{P(\arbNA\cap\arbNB)}
  \fi}
\newcommand{\feuille}[2]{\node[right,font=\footnotesize] at (#1.east) {#2};}
\newcommand{\arbval}[2]{\expandafter\def\csname arbv#1\endcsname{#2}}
\newcommand{\arbapres}[2]{\expandafter\def\csname arba#1\endcsname{#2}}
\newcommand{\arbretroisepreuves}[9]{%
  \node (O) at (0,0) {};
  \foreach \a/\ya/\na/\pa/\pos in {1/2.2/{#1}/{#4}/above,0/-2.2/{\overline{#1}}/{#5}/below}{%
    \node[evt] (N\a) at (2.9,\ya) {$\na$};
    \draw (O.center) -- node[nom,\pos]{$P(\na)=\pa$} (N\a);
    \foreach \b/\db/\nb/\pb/\posb in {1/1.1/{#2}/{#6}/above,0/-1.1/{\overline{#2}}/{#7}/below}{%
      \pgfmathsetmacro{\yb}{\ya+\db}
      \node[evt] (N\a\b) at (5.8,\yb) {$\nb$};
      \draw (N\a) -- node[nom,\posb]{$P(\nb)=\pb$} (N\a\b);
      \foreach \c/\dc/\nc/\pc/\posc in {1/0.5/{#3}/{#8}/above,0/-0.5/{\overline{#3}}/{#9}/below}{%
        \pgfmathsetmacro{\yc}{\yb+\dc}
        \node[evt] (N\a\b\c) at (8.8,\yc) {$\nc$};
        \draw (N\a\b) -- node[nom,\posc]{$P(\nc)=\pc$} (N\a\b\c);
        \node[right,font=\footnotesize] at (N\a\b\c.east)
          {$P\big((\na\,;\nb\,;\nc)\big)\ifcsname arbv\a\b\c\endcsname=\csname arbv\a\b\c\endcsname\fi
          \ifcsname arba\a\b\c\endcsname\csname arba\a\b\c\endcsname\fi$};}}}}
\newcommand{\arbretrois}[2]{\arbretroisepreuves{S}{S}{S}{#1}{#2}{#1}{#2}{#1}{#2}}

\begin{remarquebox}[Durée du chapitre]
À titre indicatif, ce chapitre est prévu pour un volume de $8$~h d'enseignement
($4$ séances de $2$~h). Une partie importante des exercices est à faire en autonomie.
\end{remarquebox}

\partiecours{Rappels : probabilités conditionnelles}

En première, on a appris à représenter une situation par un arbre pondéré et à
calculer une probabilité conditionnelle. Ces outils reviennent dans presque tous les
exercices de probabilités du bac : on les rassemble ici.

\begin{proprietebox}[Propriété 1 --- Probabilités conditionnelles et arbres]
Soit $A$ et $B$ deux événements, avec $P(A)\neq 0$.
\begin{itemize}
  \item La probabilité de $B$ \textbf{sachant} $A$ est
        $P_{A}(B)=\dfrac{P(A\cap B)}{P(A)}$ ; donc $P(A\cap B)=P(A)\times P_{A}(B)$.
  \item Dans un arbre pondéré, les branches issues d'un même nœud ont une somme
        égale à $1$, et la probabilité d'un chemin est le \textbf{produit} des
        probabilités de ses branches.
  \item \textbf{Formule des probabilités totales} : $A$ et $\overline{A}$ forment une
        partition de l'univers, donc $P(B)=P(A\cap B)+P(\overline{A}\cap B)$.
  \item $A$ et $B$ sont \textbf{indépendants} lorsque $P(A\cap B)=P(A)\times P(B)$.
\end{itemize}
\end{proprietebox}

\begin{methodebox}[Méthode 1 --- Calculer avec un arbre et la formule des probabilités totales]
\etapeexemples
Une usine fabrique des panneaux solaires ; $4\,\%$ d'entre eux présentent un
micro-défaut. Chaque panneau passe un test automatique : si le panneau a un défaut, le
test est positif dans $95\,\%$ des cas ; s'il n'en a pas, le test est positif dans
$2\,\%$ des cas. On choisit un panneau au hasard et on note $D$ : \og le panneau a un
défaut\fg{} et $T$ : \og le test est positif\fg.
\begin{questions}[label=\alph*)]
  \item Représenter la situation par un arbre pondéré.
  \item Calculer la probabilité que le test soit positif.
  \item Le test d'un panneau est positif. Calculer la probabilité que ce panneau ait un
        défaut, arrondie au millième.
\end{questions}
\begin{correction}
\cas{a)}
\begin{center}
\begin{tikzpicture}
  \arbfeuilles{0{,}038}{}{0{,}0192}{}
  \arbre{D}{T}{0{,}04}{0{,}96}{0{,}95}{0{,}05}{0{,}02}{0{,}98}
\end{tikzpicture}
\end{center}
\cas{b)}
\explique{$D$ et $\overline{D}$ forment une partition de l'univers ; d'après la formule
des probabilités totales :}{deux chemins mènent à $T$}
\begin{calculs}
P(T) &= P(D\cap T)+P(\overline{D}\cap T) &\qquad& \rem{probabilités totales}\\
P(T) &= P(D)\times P_{D}(T)+P(\overline{D})\times P_{\overline{D}}(T) && \rem{règle du chemin}\\
P(T) &= 0{,}04\times 0{,}95+0{,}96\times 0{,}02 &&\\
P(T) &= 0{,}038+0{,}0192 &&\\
P(T) &= \res{0{,}0572} &&
\end{calculs}
\explique{La probabilité que le test soit positif est égale à $0{,}0572$.}{on conclut
par une phrase}
\cas{c)}
\explique{On cherche $P_{T}(D)$.}{\og sachant que le test est positif\fg}
\begin{calculs}
P_{T}(D) &= \frac{P(D\cap T)}{P(T)} &\qquad& \rem{définition}\\
P_{T}(D) &= \frac{0{,}038}{0{,}0572} &&\\
P_{T}(D) &\approx \res{0{,}664} &&
\end{calculs}
\explique{Sachant que son test est positif, la probabilité qu'un panneau ait un défaut
est d'environ $0{,}664$.}{le test se trompe encore une fois sur trois}
\end{correction}

\etape{À vous de jouer}
Une usine fabrique des batteries de vélo électrique ; $5\,\%$ d'entre elles présentent
un défaut. Si la batterie a un défaut, le test de contrôle est positif dans $90\,\%$ des
cas ; sinon, il est positif dans $3\,\%$ des cas. On choisit une batterie au hasard ;
$D$ : \og la batterie a un défaut\fg{} et $T$ : \og le test est positif\fg.
\begin{questions}[label=\alph*)]
  \item Représenter la situation par un arbre pondéré.
  \item Calculer la probabilité que le test soit positif.
  \item Le test d'une batterie est positif. Calculer la probabilité que cette batterie
        ait un défaut, arrondie au millième.
\end{questions}
\reponse{b) $P(T)=0{,}045+0{,}0285=0{,}0735$.\; c)~$P_{T}(D)=\frac{0{,}045}{0{,}0735}
\approx 0{,}612$.}
\end{methodebox}

\afaire{exercices 1 à 4}

\partiecours{Rappels : variable aléatoire, loi, espérance}

\begin{definitionbox}[Définition 1 --- Variable aléatoire, loi, espérance]
Une \textbf{variable aléatoire} $X$ associe un nombre réel à chaque issue d'une
expérience aléatoire (un gain, un nombre de succès\dots). On note $\{X=a\}$ l'événement
\og $X$ prend la valeur $a$\fg.

\smallskip
Si $X$ prend les valeurs $x_{1}$, $x_{2}$, \dots, $x_{n}$, sa \textbf{loi de probabilité}
donne la probabilité $p_{i}=P(X=x_{i})$ de chaque valeur ; la somme des $p_{i}$ vaut $1$.
\begin{itemize}
  \item \textbf{Espérance} : $E(X)=p_{1}x_{1}+p_{2}x_{2}+\dots+p_{n}x_{n}$ ; c'est la
        moyenne des valeurs de $X$ sur un grand nombre de répétitions.
  \item \textbf{Variance} : $V(X)=p_{1}\big(x_{1}-E(X)\big)^{2}+\dots
        +p_{n}\big(x_{n}-E(X)\big)^{2}$ ; \textbf{écart-type} : $\sigma(X)=\sqrt{V(X)}$,
        qui mesure la dispersion des valeurs autour de $E(X)$.
\end{itemize}
Pour un gain, le jeu est \textbf{favorable} au joueur si $E(X)>0$, \textbf{équitable} si
$E(X)=0$, \textbf{défavorable} si $E(X)<0$.
\end{definitionbox}

\begin{methodebox}[Méthode 2 --- Déterminer une loi et calculer une espérance]
\etapeexemples
À une fête de quartier, on mise $3$\,€ et on lance un dé équilibré à six faces : le
$6$ fait gagner $10$\,€, le $4$ et le $5$ font gagner $2$\,€, les autres faces ne font
rien gagner. On note $X$ le gain algébrique du joueur : ce qu'il reçoit, moins sa mise.
\begin{questions}[label=\alph*)]
  \item Déterminer la loi de probabilité de $X$.
  \item Calculer $E(X)$ et interpréter le résultat.
  \item Quelle mise rendrait le jeu équitable ?
\end{questions}
\begin{correction}
\cas{a)}
\explique{$X$ prend les valeurs $10-3=7$ (face $6$), $2-3=-1$ (faces $4$ et $5$) et
$0-3=-3$ (faces $1$, $2$, $3$).}{le dé est équilibré : $\frac{1}{6}$ par face}
\begin{center}
\renewcommand{\arraystretch}{2.4}\setlength{\tabcolsep}{10pt}
\begin{tabular}{|c|c|c|c|}
\hline
$x_{i}$ & $-3$ & $-1$ & $7$\\\hline
$P(X=x_{i})$ & $\dfrac{3}{6}$ & $\dfrac{2}{6}$ & $\dfrac{1}{6}$\\\hline
\end{tabular}
\end{center}
\cas{b)}
\begin{calculs}
E(X) &= (-3)\times\frac{3}{6}+(-1)\times\frac{2}{6}+7\times\frac{1}{6} &\qquad& \rem{somme
des $p_{i}x_{i}$}\\
E(X) &= \frac{-9-2+7}{6} &&\\
E(X) &= -\frac{2\times\barre{2}}{3\times\barre{2}} = \res{-\frac{2}{3}} && \rem{on
simplifie par $2$}
\end{calculs}
\explique{$E(X)<0$ : le jeu est défavorable au joueur ; sur un grand nombre de parties,
il perd en moyenne environ $0{,}67$\,€ par partie.}{interprétation}
\cas{c)}
\explique{Avec une mise de $m$ euros, le gain algébrique vaut $10-m$, $2-m$ ou $-m$, et
son espérance est $\frac{1}{6}(10-m)+\frac{2}{6}(2-m)+\frac{3}{6}(-m)=\frac{14-6m}{6}$.
Le jeu est équitable lorsque cette espérance est nulle :}{même loi, autres valeurs}
\begin{calculs}
\frac{14-6m}{6} &= 0 &\qquad& \rem{espérance nulle}\\
\iff 14-6m &= 0 && \rem{on multiplie par $6$}\\
\iff m &= \frac{14}{6} = \frac{7}{3} &&
\end{calculs}
\explique{Le jeu est équitable pour une mise de \res{$\frac{7}{3}$\,€}, environ
$2{,}33$\,€.}{on conclut par une phrase}
\end{correction}

\etape{À vous de jouer}
On mise $2$\,€ et on lance un dé équilibré à six faces : le $1$ fait gagner $6$\,€, le
$2$ et le $3$ font gagner $2$\,€, les autres faces ne font rien gagner. On note $X$ le
gain algébrique du joueur.
\begin{questions}[label=\alph*)]
  \item Déterminer la loi de probabilité de $X$.
  \item Calculer $E(X)$ et interpréter le résultat.
  \item Quelle mise rendrait le jeu équitable ?
\end{questions}
\reponse{a) $P(X=-2)=\frac{3}{6}$ ; $P(X=0)=\frac{2}{6}$ ; $P(X=4)=\frac{1}{6}$.\;
b)~$E(X)=-\frac{1}{3}$ : jeu défavorable.\; c)~$m=\frac{5}{3}$\,€, environ $1{,}67$\,€.}
\end{methodebox}

\afaire{exercices 5 à 7}

\partiecours{Succession d'épreuves indépendantes}

\begin{definitionbox}[Définition 2 --- Épreuves indépendantes]
Des épreuves qui se succèdent sont \textbf{indépendantes} lorsque l'issue de chacune ne
dépend pas des issues des précédentes.

\smallskip
Si les épreuves ont pour univers $E_{1}$, $E_{2}$, \dots, $E_{n}$, l'univers de leur
succession est le produit cartésien $E_{1}\times E_{2}\times\dots\times E_{n}$ (chapitre
3) : une issue est un $n$-uplet $(x_{1}\,;x_{2}\,;\dots\,;x_{n})$, avec $x_{i}\in E_{i}$.

\smallskip
Exemple : on lance une pièce, puis un dé ; le résultat de la pièce n'influe pas sur
celui du dé. L'univers est $\{P\,;F\}\times\{1\,;2\,;3\,;4\,;5\,;6\}$, qui a
$2\times 6=12$ issues, comme $(F\,;4)$.
\end{definitionbox}

\begin{proprietebox}[Propriété 2 --- Probabilité d'une issue (admise)]
Dans une succession d'épreuves indépendantes, la probabilité d'une issue
$(x_{1}\,;x_{2}\,;\dots\,;x_{n})$ est le \textbf{produit} des probabilités de ses
composantes :
\[
P\big((x_{1}\,;x_{2}\,;\dots\,;x_{n})\big)=P(x_{1})\times P(x_{2})\times\dots\times P(x_{n}).
\]
Sur l'arbre, chaque branche porte la probabilité de l'issue de son épreuve, la même à
chaque niveau quel que soit le chemin suivi : c'est la règle du chemin.
\end{proprietebox}

\begin{methodebox}[Méthode 3 --- Calculer une probabilité dans une succession d'épreuves indépendantes]
\etapeexemples
Pour aller au lycée à vélo, Inès passe deux feux tricolores, puis un passage à niveau,
qui fonctionnent indépendamment. Le premier feu est vert avec une probabilité de
$0{,}6$, le second avec une probabilité de $0{,}5$, et la barrière est levée avec une
probabilité de $0{,}9$. On note $V_{1}$ : \og le premier feu est vert\fg, $V_{2}$ :
\og le second feu est vert\fg{} et $L$ : \og la barrière est levée\fg.
\begin{questions}[label=\alph*)]
  \item Représenter la situation par un arbre pondéré.
  \item Calculer la probabilité qu'Inès passe sans s'arrêter.
  \item Calculer la probabilité qu'elle s'arrête exactement une fois.
  \item Calculer la probabilité qu'elle s'arrête au moins une fois.
\end{questions}
\begin{correction}
\cas{a)}
\explique{Les trois épreuves sont indépendantes : chaque branche porte la probabilité
de son épreuve, quel que soit le chemin ; on écrit $P(V_{2})$ plutôt que
$P_{V_{1}}(V_{2})$.}{$P(\overline{V_{1}})=1-0{,}6=0{,}4$\dots}
\begin{center}
\begin{tikzpicture}
  \arbval{111}{0{,}27}\arbval{011}{0{,}18}\arbval{101}{0{,}27}\arbval{110}{0{,}03}
  \arbretroisepreuves{V_{1}}{V_{2}}{L}{0{,}6}{0{,}4}{0{,}5}{0{,}5}{0{,}9}{0{,}1}
\end{tikzpicture}
\end{center}
\cas{b)}
\explique{Inès passe sans s'arrêter pour l'issue $(V_{1}\,;V_{2}\,;L)$.}{un seul chemin}
\begin{calculs}
P\big((V_{1}\,;V_{2}\,;L)\big) &= 0{,}6\times 0{,}5\times 0{,}9 &\qquad&
\rem{épreuves indépendantes : produit}\\
P\big((V_{1}\,;V_{2}\,;L)\big) &= \res{0{,}27} &&
\end{calculs}
\cas{c)}
\explique{Trois issues comptent exactement un arrêt : $(\overline{V_{1}}\,;V_{2}\,;L)$,
$(V_{1}\,;\overline{V_{2}}\,;L)$ et $(V_{1}\,;V_{2}\,;\overline{L})$. On note $A$ cet
événement.}{l'arrêt a lieu au 1\textsuperscript{er} feu, au 2\textsuperscript{e} ou à la
barrière}
\begin{calculs}
P(A) &= 0{,}4\times 0{,}5\times 0{,}9+0{,}6\times 0{,}5\times 0{,}9+0{,}6\times 0{,}5
\times 0{,}1 &\qquad& \rem{somme des issues}\\
P(A) &= 0{,}18+0{,}27+0{,}03 &&\\
P(A) &= \res{0{,}48} &&
\end{calculs}
\cas{d)}
\explique{\og S'arrêter au moins une fois\fg{} est l'événement contraire de \og passer
sans s'arrêter\fg, dont la probabilité vaut $0{,}27$ ; on note $M$ cet événement.}{plus
court que d'additionner sept issues}
\begin{calculs}
P(M) &= 1-0{,}27 &\qquad& \rem{événement contraire}\\
P(M) &= \res{0{,}73} &&
\end{calculs}
\end{correction}

\etape{À vous de jouer}
Pour aller au lycée à vélo, Malo passe deux feux tricolores, puis un passage à niveau,
qui fonctionnent indépendamment. Le premier feu est vert avec une probabilité de
$0{,}7$, le second avec une probabilité de $0{,}4$, et la barrière est levée avec une
probabilité de $0{,}8$. On note $V_{1}$, $V_{2}$ et $L$ les mêmes événements.
\begin{questions}[label=\alph*)]
  \item Représenter la situation par un arbre pondéré.
  \item Calculer la probabilité que Malo passe sans s'arrêter.
  \item Calculer la probabilité qu'il s'arrête exactement une fois.
  \item Calculer la probabilité qu'il s'arrête au moins une fois.
\end{questions}
\reponse{b) $0{,}7\times 0{,}4\times 0{,}8=0{,}224$.\; c)~$0{,}096+0{,}336+0{,}056=0{,}488$.\;
d)~$1-0{,}224=0{,}776$.}
\end{methodebox}

\afaire{exercices 8 à 11}

\partiecours{Épreuve, loi et schéma de Bernoulli}

\begin{definitionbox}[Définition 3 --- Épreuve et loi de Bernoulli]
Une \textbf{épreuve de Bernoulli} est une expérience aléatoire à deux issues : le
\textbf{succès} $S$, de probabilité $p$, et l'\textbf{échec} $\overline{S}$, de
probabilité $1-p$.

\smallskip
La variable aléatoire $X$ qui vaut $1$ en cas de succès et $0$ en cas d'échec suit la
\textbf{loi de Bernoulli de paramètre $p$} : $P(X=1)=p$ et $P(X=0)=1-p$.

\smallskip
Exemple : on tire une boule d'une urne de $10$ boules, dont $3$ rouges ; le succès
\og la boule est rouge\fg{} a pour probabilité $p=0{,}3$. Une expérience à plusieurs
issues devient une épreuve de Bernoulli dès qu'on choisit le succès.
\end{definitionbox}

\begin{proprietebox}[Propriété 3 --- Espérance et variance d'une loi de Bernoulli]
Si $X$ suit la loi de Bernoulli de paramètre $p$, alors $E(X)=p$ et $V(X)=p(1-p)$.
\end{proprietebox}

\textbf{\color{coulMeth}Démonstration.} $X$ prend les valeurs $0$ et $1$, avec les
probabilités $1-p$ et $p$.
\begin{calculs}
E(X) &= (1-p)\times 0+p\times 1 = p &\qquad& \rem{somme des $p_{i}x_{i}$}\\[3pt]
V(X) &= (1-p)\times(0-p)^{2}+p\times(1-p)^{2} && \rem{définition, avec $E(X)=p$}\\
V(X) &= p(1-p)\times p+p(1-p)\times(1-p) && \rem{$(1-p)p^{2}=p(1-p)\times p$}\\
V(X) &= p(1-p)\big(p+(1-p)\big) && \rem{on factorise par $p(1-p)$}\\
V(X) &= p(1-p) &&
\end{calculs}

\begin{definitionbox}[Définition 4 --- Schéma de Bernoulli]
Un \textbf{schéma de Bernoulli} de paramètres $n$ et $p$ est la répétition de $n$
épreuves de Bernoulli \textbf{identiques} et \textbf{indépendantes}, de même probabilité
de succès $p$. Son univers est $\{S\,;\overline{S}\}^{n}$ ; sur l'arbre, chaque branche
$S$ porte $p$ et chaque branche $\overline{S}$ porte $1-p$.
\end{definitionbox}

\begin{remarquebox}[Tirage avec remise, tirage sans remise]
Des tirages \textbf{avec remise} sont identiques et indépendants : ils forment un schéma
de Bernoulli. Des tirages \textbf{sans remise} n'en forment pas un, car la composition
de l'urne change à chaque tirage. Quand on prélève quelques objets dans un très grand
stock (une production, une population), les énoncés de bac disent que \og ce choix est
assimilé à un tirage avec remise\fg{} : on peut alors parler de schéma de Bernoulli.
\end{remarquebox}

\begin{methodebox}[Méthode 4 --- Justifier qu'une situation est un schéma de Bernoulli]
\etapeexemples
Le premier étage d'une fusée réutilisable revient se poser sur sa barge avec une
probabilité de $0{,}9$. On étudie trois lancements ; les atterrissages sont supposés
indépendants les uns des autres.
\begin{questions}[label=\alph*)]
  \item Décrire l'épreuve de Bernoulli associée à un lancement.
  \item Justifier que la situation est un schéma de Bernoulli et donner ses paramètres.
  \item Représenter ce schéma par un arbre pondéré.
  \item Calculer la probabilité qu'au moins un atterrissage échoue.
\end{questions}
\begin{correction}
\cas{a)}
\explique{Un lancement est une épreuve à deux issues : le succès $S$ \og le premier étage
se pose sur sa barge\fg, de probabilité \res{$p=0{,}9$}, et l'échec $\overline{S}$, de
probabilité $1-0{,}9=0{,}1$.}{on nomme le succès et on donne sa probabilité}
\cas{b)}
\explique{On répète trois fois cette épreuve de Bernoulli, de façon identique et
indépendante : c'est un schéma de Bernoulli de paramètres \res{$n=3$ et
$p=0{,}9$}.}{\og identique\fg{} : même $p$ ; \og indépendante\fg{} : l'énoncé le dit}
\cas{c)}
\begin{center}
\begin{tikzpicture}
  \arbval{111}{0{,}9^{3}}
  \arbretrois{0{,}9}{0{,}1}
\end{tikzpicture}
\end{center}
\cas{d)}
\explique{\og Au moins un échec\fg{} est l'événement contraire de \og trois succès\fg,
l'issue $(S\,;S\,;S)$ ; on note $E$ cet événement.}{un seul chemin à retirer}
\begin{calculs}
P(E) &= 1-0{,}9\times 0{,}9\times 0{,}9 &\qquad& \rem{événement contraire, règle du
chemin}\\
P(E) &= 1-0{,}729 &&\\
P(E) &= \res{0{,}271} &&
\end{calculs}
\end{correction}

\etape{À vous de jouer}
Un panneau solaire est sans défaut avec une probabilité de $0{,}97$. On prélève trois
panneaux dans un stock assez grand pour assimiler ce choix à un tirage avec remise.
\begin{questions}[label=\alph*)]
  \item Décrire l'épreuve de Bernoulli associée à un panneau.
  \item Justifier que la situation est un schéma de Bernoulli et donner ses paramètres.
  \item Représenter ce schéma par un arbre pondéré.
  \item Calculer la probabilité qu'au moins un des trois panneaux ait un défaut.
\end{questions}
\reponse{a) $S$ : \og le panneau est sans défaut\fg, $p=0{,}97$.\; b)~$n=3$ et
$p=0{,}97$.\; d)~$1-0{,}97^{3}\approx 0{,}087$.}
\end{methodebox}

\afaire{exercices 12 à 15}

\partiecours{Loi binomiale}

\textbf{Pour commencer.} Un sac contient $3$ jetons rouges et $7$ jetons jaunes. On tire
un jeton, on note sa couleur, on le remet, et on recommence : trois tirages en tout. Le
succès $S$ est \og le jeton est rouge\fg, de probabilité $0{,}3$. On note $X$ le nombre
de jetons rouges obtenus.

\begin{center}
\begin{tikzpicture}
  \arbval{111}{0{,}3^{3}}\arbapres{111}{\qquad X=3}
  \arbval{110}{0{,}3^{2}\times 0{,}7}\arbapres{110}{\qquad X=2}
  \arbval{101}{0{,}3^{2}\times 0{,}7}\arbapres{101}{\qquad X=2}
  \arbval{100}{0{,}3\times 0{,}7^{2}}\arbapres{100}{\qquad X=1}
  \arbval{011}{0{,}3^{2}\times 0{,}7}\arbapres{011}{\qquad X=2}
  \arbval{010}{0{,}3\times 0{,}7^{2}}\arbapres{010}{\qquad X=1}
  \arbval{001}{0{,}3\times 0{,}7^{2}}\arbapres{001}{\qquad X=1}
  \arbval{000}{0{,}7^{3}}\arbapres{000}{\qquad X=0}
  \arbretrois{0{,}3}{0{,}7}
\end{tikzpicture}
\end{center}
Trois chemins mènent à $X=2$ ; chacun compte deux succès et un échec, donc a pour
probabilité $0{,}3^{2}\times 0{,}7$.

\smallskip
${P(X=2)=3\times 0{,}3^{2}\times 0{,}7}$, soit ${P(X=2)=0{,}189}$.

\smallskip
De même, un seul chemin mène à $X=3$ : ${P(X=3)=0{,}3^{3}}$, soit $0{,}027$.

\smallskip
\begin{minipage}[c]{0.56\linewidth}\raggedright
On obtient la loi de $X$ :
\begin{center}
\renewcommand{\arraystretch}{1.4}\setlength{\tabcolsep}{8pt}
\begin{tabular}{|c|c|c|c|c|}
\hline
$k$ & $0$ & $1$ & $2$ & $3$\\\hline
$P(X=k)$ & $0{,}343$ & $0{,}441$ & $0{,}189$ & $0{,}027$\\\hline
\end{tabular}
\end{center}
Le nombre de chemins qui comptent $k$ succès parmi $3$ épreuves est $\binom{3}{k}$
(chapitre 3, partie VIII) : $1$, $3$, $3$ et $1$.
\end{minipage}\hfill
\begin{minipage}[c]{0.4\linewidth}\centering
\begin{tikzpicture}[x=0.9cm,y=5cm]
  \draw[->] (-0.6,0) -- (3.7,0) node[right,font=\footnotesize]{$k$};
  \draw[->] (-0.6,0) -- (-0.6,0.52);
  \foreach \k/\v in {0/0.343,1/0.441,2/0.189,3/0.027}{
    \fill[coulMeth!60] (\k-0.18,0) rectangle (\k+0.18,\v);
    \node[below,font=\footnotesize] at (\k,0) {$\k$};}
  \foreach \y/\t in {0.1/0{,}1,0.2/0{,}2,0.3/0{,}3,0.4/0{,}4}{
    \draw (-0.65,\y) -- (-0.55,\y) node[left=2pt,font=\footnotesize]{$\t$};}
\end{tikzpicture}
\end{minipage}

\begin{definitionbox}[Définition 5 --- Loi binomiale]
On considère un schéma de Bernoulli de paramètres $n$ et $p$, et $X$ la variable
aléatoire qui compte le \textbf{nombre de succès}. La loi de $X$ est la \textbf{loi
binomiale} de paramètres $n$ et $p$, notée $\mathcal{B}(n\,;p)$.

\smallskip
$X$ prend les valeurs entières de $0$ (aucun succès) à $n$ (que des succès).
\end{definitionbox}

\begin{proprietebox}[Propriété 4 --- Probabilité de $k$ succès]
Si $X$ suit la loi binomiale $\mathcal{B}(n\,;p)$, alors, pour tout entier $k$ tel que
$0\leq k\leq n$ :
\[
P(X=k)=\binom{n}{k}\,p^{k}\,(1-p)^{n-k}.
\]
\end{proprietebox}

\textbf{\color{coulMeth}Démonstration} (exigible).
Soit $k$ un entier tel que $0\leq k\leq n$.
\begin{itemize}
  \item Un chemin de l'arbre qui réalise $k$ succès compte $k$ branches $S$, de
        probabilité $p$, et $n-k$ branches $\overline{S}$, de probabilité $1-p$. Les
        épreuves sont indépendantes, donc ce chemin a pour probabilité
        $p^{k}(1-p)^{n-k}$.
  \item Un tel chemin correspond au choix des $k$ épreuves réussies parmi les $n$ : il
        y a $\binom{n}{k}$ chemins qui réalisent $k$ succès (chapitre 3).
  \item Ces chemins sont des issues distinctes, donc on additionne leurs probabilités,
        toutes égales.
\end{itemize}
Donc $P(X=k)=\dbinom{n}{k}\,p^{k}\,(1-p)^{n-k}$.

\begin{methodebox}[Méthode 5 --- Reconnaître une loi binomiale et calculer la probabilité de $k$ succès]
\etapeexemples
Dans une ville, $35\,\%$ des habitants vont au travail à vélo. On interroge au hasard
$12$ habitants ; la population est assez grande pour assimiler ce choix à un tirage avec
remise. On note $X$ le nombre d'habitants interrogés qui vont au travail à vélo.
\begin{questions}[label=\alph*)]
  \item Justifier que $X$ suit une loi binomiale et préciser ses paramètres.
  \item Calculer la probabilité qu'exactement $4$ des habitants interrogés aillent au
        travail à vélo. Donner l'expression exacte, puis une valeur arrondie au millième.
  \item Calculer la probabilité qu'aucun n'y aille à vélo, arrondie au millième.
\end{questions}
\begin{correction}
\cas{a)}
\explique{Interroger un habitant est une épreuve de Bernoulli, de succès \og l'habitant va
au travail à vélo\fg, de probabilité $p=0{,}35$.}{l'épreuve et son succès}
\explique{Le choix étant assimilé à un tirage avec remise, on répète $12$ fois cette
épreuve de façon identique et indépendante.}{le schéma de Bernoulli}
\explique{$X$ compte le nombre de succès, donc $X$ suit la loi binomiale
\res{$\mathcal{B}(12\,;0{,}35)$}.}{la conclusion, avec $n$ et $p$}
\cas{b)}
\begin{calculs}
P(X=4) &= \binom{12}{4}\times 0{,}35^{4}\times(1-0{,}35)^{12-4} &\qquad& \rem{propriété 4,
avec $k=4$}\\
P(X=4) &= \res{495\times 0{,}35^{4}\times 0{,}65^{8}} && \rem{$\binom{12}{4}=495$, à la
calculatrice}\\
P(X=4) &\approx \res{0{,}237} &&
\end{calculs}
\explique{La probabilité qu'exactement $4$ des habitants interrogés aillent au travail à
vélo est d'environ $0{,}237$.}{on conclut par une phrase}
\cas{c)}
\begin{calculs}
P(X=0) &= \binom{12}{0}\times 0{,}35^{0}\times 0{,}65^{12} &\qquad& \rem{$k=0$}\\
P(X=0) &= 0{,}65^{12} && \rem{$\binom{12}{0}=1$ et $0{,}35^{0}=1$}\\
P(X=0) &\approx \res{0{,}006} &&
\end{calculs}
\end{correction}

\etape{À vous de jouer}
Dans un collège, $25\,\%$ des élèves jouent d'un instrument de musique. On interroge au
hasard $10$ élèves, en assimilant ce choix à un tirage avec remise. On note $X$ le
nombre d'élèves interrogés qui jouent d'un instrument.
\begin{questions}[label=\alph*)]
  \item Justifier que $X$ suit une loi binomiale et préciser ses paramètres.
  \item Calculer la probabilité qu'exactement $3$ des élèves interrogés jouent d'un
        instrument. Donner l'expression exacte, puis une valeur arrondie au millième.
  \item Calculer la probabilité qu'aucun ne joue d'un instrument, arrondie au millième.
\end{questions}
\reponse{a) $\mathcal{B}(10\,;0{,}25)$.\; b)~$120\times 0{,}25^{3}\times 0{,}75^{7}\approx
0{,}250$.\; c)~$0{,}75^{10}\approx 0{,}056$.}
\end{methodebox}

\afaire{exercices 16 à 20}

\partiecours{Probabilités cumulées à la calculatrice}

Les questions de bac portent souvent sur \og au plus\fg, \og au moins\fg{} ou \og entre
\dots{} et \dots\fg{} : il faudrait additionner beaucoup de probabilités $P(X=k)$. La
calculatrice donne directement $P(X\leq k)$, la somme $P(X=0)+P(X=1)+\dots+P(X=k)$.

\begin{proprietebox}[Propriété 5 --- Traduire un énoncé]
Soit $X$ une variable aléatoire à valeurs entières et $k$ un entier.
\begin{center}
\renewcommand{\arraystretch}{1.5}\setlength{\tabcolsep}{8pt}
\begin{tabular}{|l|l|l|}
\hline
\og au plus $k$\fg{} & $X\leq k$ & $P(X\leq k)$, directement\\\hline
\og moins de $k$\fg{} & $X<k$, soit $X\leq k-1$ & $P(X\leq k-1)$\\\hline
\og au moins $k$\fg{} & $X\geq k$ & $1-P(X\leq k-1)$, par l'événement contraire\\\hline
\og plus de $k$\fg{} & $X>k$ & $1-P(X\leq k)$\\\hline
\og entre $a$ et $b$\fg{} (inclus) & $a\leq X\leq b$ & $P(X\leq b)-P(X\leq a-1)$\\\hline
\end{tabular}
\end{center}
Le contraire de $\{X\geq k\}$ est $\{X\leq k-1\}$, pas $\{X\leq k\}$ : la valeur $k$
appartient à $\{X\geq k\}$.
\end{proprietebox}

\begin{remarquebox}[Loi binomiale à la calculatrice]
\begin{itemize}
  \item \textbf{NumWorks} : application Probabilités, Binomiale, saisir $n$ et $p$, puis
        choisir la forme $P(X\leq k)$, $P(X=k)$ ou $P(a\leq X\leq b)$.
  \item \textbf{TI} : \texttt{2nde} \texttt{var} (distrib) ; \texttt{binomFdp(n,p,k)}
        donne $P(X=k)$ et \texttt{binomFRép(n,p,k)} donne $P(X\leq k)$.
  \item \textbf{Casio} : menu STAT, DIST, BINOMIAL ; \texttt{Bpd} donne $P(X=k)$ et
        \texttt{Bcd} donne $P(a\leq X\leq b)$ (avec $a=0$ : $P(X\leq b)$).
\end{itemize}
Au bac, on écrit la probabilité cherchée, la traduction, puis \og à la calculatrice,
$P(X\leq 25)\approx 0{,}683$\fg.
\end{remarquebox}

\begin{methodebox}[Méthode 6 --- Calculer une probabilité \og au plus\fg, \og au moins\fg, \og entre\fg]
\etapeexemples
Une société d'autopartage dispose de $40$ voitures électriques. Un jour donné, chaque
voiture est réservée avec une probabilité de $0{,}6$, indépendamment des autres. On
admet que le nombre $X$ de voitures réservées suit la loi binomiale
$\mathcal{B}(40\,;0{,}6)$. Calculer, arrondie au millième, la probabilité que :
\begin{questions}[label=\alph*)]
  \item au plus $25$ voitures soient réservées ;
  \item au moins $20$ voitures soient réservées ;
  \item entre $20$ et $28$ voitures (incluses) soient réservées ;
  \item plus de $30$ voitures soient réservées.
\end{questions}
\begin{correction}
\cas{a)}
\begin{calculs}
P(X\leq 25) &\approx \res{0{,}683} &\qquad& \rem{à la calculatrice, directement}
\end{calculs}
\cas{b)}
\explique{L'événement contraire de $\{X\geq 20\}$ est $\{X\leq 19\}$.}{$20$ appartient à
$\{X\geq 20\}$}
\begin{calculs}
P(X\geq 20) &= 1-P(X\leq 19) &\qquad& \rem{événement contraire}\\
P(X\geq 20) &\approx \res{0{,}926} && \rem{à la calculatrice}
\end{calculs}
\cas{c)}
\begin{calculs}
P(20\leq X\leq 28) &= P(X\leq 28)-P(X\leq 19) &\qquad& \rem{on retire les valeurs de $0$ à
$19$}\\
P(20\leq X\leq 28) &\approx \res{0{,}855} && \rem{à la calculatrice}
\end{calculs}
\cas{d)}
\explique{\og Plus de $30$\fg{} se traduit par $X>30$, soit $X\geq 31$ ; son contraire est
$\{X\leq 30\}$.}{$30$ n'est pas compris}
\begin{calculs}
P(X>30) &= 1-P(X\leq 30) &\qquad& \rem{événement contraire}\\
P(X>30) &\approx \res{0{,}016} && \rem{à la calculatrice}
\end{calculs}
\end{correction}

\etape{À vous de jouer}
Un club d'astronomie compte $30$ membres. Chaque membre vient à la soirée d'observation
avec une probabilité de $0{,}45$, indépendamment des autres. On admet que le nombre $X$
de membres présents suit la loi binomiale $\mathcal{B}(30\,;0{,}45)$. Calculer, arrondie
au millième, la probabilité que :
\begin{questions}[label=\alph*)]
  \item au plus $15$ membres soient présents ;
  \item au moins $11$ membres soient présents ;
  \item entre $11$ et $18$ membres (inclus) soient présents ;
  \item plus de $20$ membres soient présents.
\end{questions}
\reponse{a) $0{,}769$.\; b)~$1-P(X\leq 10)\approx 0{,}865$.\; c)~$P(X\leq 18)-P(X\leq 10)
\approx 0{,}832$.\; d)~$1-P(X\leq 20)\approx 0{,}005$.}
\end{methodebox}

\afaire{exercices 21 à 25}

\partiecours{Problèmes de seuil}

Deux questions reviennent au bac : combien de répétitions faut-il pour être presque sûr
d'obtenir au moins un succès ? Quel nombre $k$ ne sera dépassé qu'avec une faible
probabilité ? Dans les deux cas, on cherche le \textbf{plus petit} entier qui convient.

\begin{proprietebox}[Propriété 6 --- Au moins un succès]
Si $X$ suit la loi binomiale $\mathcal{B}(n\,;p)$, alors $P(X=0)=(1-p)^{n}$ et, par
l'événement contraire,
\[
P(X\geq 1)=1-(1-p)^{n}.
\]
\end{proprietebox}

\begin{methodebox}[Méthode 7 --- Trouver le nombre de répétitions pour obtenir au moins un succès]
\etapeexemples
Dans une réserve naturelle, un piège photographique a, chaque nuit, une probabilité de
$0{,}12$ de photographier un lynx, indépendamment des autres nuits. Soit $n\in\Ns$ ; on
note $X_{n}$ le nombre de nuits, sur $n$ nuits, où un lynx est photographié. On admet
que $X_{n}$ suit la loi binomiale $\mathcal{B}(n\,;0{,}12)$.
\begin{questions}[label=\alph*)]
  \item Exprimer en fonction de $n$ la probabilité $p_{n}$ qu'un lynx soit photographié au
        moins une fois en $n$ nuits.
  \item Déterminer le plus petit nombre de nuits $n$ pour lequel $p_{n}\geq 0{,}95$.
\end{questions}
\begin{correction}
\cas{a)}
\explique{L'événement \og au moins une photo\fg{} est le contraire de $\{X_{n}=0\}$, et
$P(X_{n}=0)=\binom{n}{0}\times 0{,}12^{0}\times 0{,}88^{n}=0{,}88^{n}$.}{propriété 4,
avec $k=0$}
\begin{calculs}
p_{n} &= P(X_{n}\geq 1) &\qquad& \rem{\og au moins une fois\fg}\\
p_{n} &= 1-P(X_{n}=0) && \rem{événement contraire}\\
p_{n} &= \res{1-0{,}88^{n}} &&
\end{calculs}
\cas{b)}
\begin{calculs}
p_{n} &\geq 0{,}95 &\qquad& \rem{la condition cherchée}\\
1-0{,}88^{n} &\geq 0{,}95 && \rem{on remplace $p_{n}$}\\
\iff -0{,}88^{n} &\geq -0{,}05 && \rem{on retire $1$}\\
\iff 0{,}88^{n} &\leq 0{,}05 && \rem{on multiplie par $-1<0$ : l'ordre change}
\end{calculs}
\explique{À la calculatrice (table des valeurs de $0{,}88^{n}$) : $0{,}88^{23}\approx
0{,}0529$ et $0{,}88^{24}\approx 0{,}0465$.}{$0{,}88^{23}>0{,}05\geq 0{,}88^{24}$}
\explique{La suite $(0{,}88^{n})$ est géométrique de raison $0{,}88$, avec
$0<0{,}88<1$ : elle est décroissante. Pour $n\leq 23$, $0{,}88^{n}\geq 0{,}88^{23}>0{,}05$ ;
pour $n\geq 24$, $0{,}88^{n}\leq 0{,}88^{24}\leq 0{,}05$.}{une fois le seuil franchi, on
ne revient pas en arrière}
\explique{Le plus petit nombre de nuits est donc \res{$n=24$}.}{on conclut par une
phrase}
\end{correction}

\begin{remarquebox}[Avec le logarithme]
Au chapitre 9, la fonction logarithme permettra de résoudre $0{,}88^{n}\leq 0{,}05$ par
le calcul, comme dans les corrigés de bac ; la méthode ci-dessus reste valable.
\end{remarquebox}

\etape{À vous de jouer}
Chaque nuit, un astronome amateur a une probabilité de $0{,}15$ de photographier un
météore brillant, indépendamment des autres nuits. Soit $n\in\Ns$ ; on note $X_{n}$ le
nombre de nuits, sur $n$ nuits, où il en photographie un. On admet que $X_{n}$ suit la
loi binomiale $\mathcal{B}(n\,;0{,}15)$.
\begin{questions}[label=\alph*)]
  \item Exprimer en fonction de $n$ la probabilité $p_{n}$ qu'il photographie un météore
        brillant au moins une fois en $n$ nuits.
  \item Déterminer le plus petit nombre de nuits $n$ pour lequel $p_{n}\geq 0{,}99$.
\end{questions}
\reponse{a) $p_{n}=1-0{,}85^{n}$.\; b)~$0{,}85^{n}\leq 0{,}01$ ; $0{,}85^{28}\approx 0{,}0106$
et $0{,}85^{29}\approx 0{,}0090$ : $n=29$.}
\end{methodebox}

\begin{methodebox}[Méthode 8 --- Trouver le plus petit entier $k$ tel que $P(X>k)\leq\alpha$]
\etapeexemples
Un musée propose une visite avec audioguide ; $200$ personnes ont réservé. Chacune vient
avec une probabilité de $0{,}85$, indépendamment des autres. On admet que le nombre $X$
de visiteurs présents suit la loi binomiale $\mathcal{B}(200\,;0{,}85)$. Le musée prépare
$k$ audioguides. Déterminer le plus petit nombre $k$ d'audioguides pour que la
probabilité d'en manquer soit inférieure ou égale à $0{,}05$.
\begin{correction}
\cas{Résolution}\par\nopagebreak
\explique{On manque d'audioguides lorsque plus de $k$ visiteurs viennent, soit
$X>k$.}{on traduit l'énoncé}
\begin{calculs}
P(X>k) &\leq 0{,}05 &\qquad& \rem{la condition cherchée}\\
\iff 1-P(X\leq k) &\leq 0{,}05 && \rem{événement contraire}\\
\iff -P(X\leq k) &\leq -0{,}95 && \rem{on retire $1$}\\
\iff P(X\leq k) &\geq 0{,}95 && \rem{on multiplie par $-1<0$ : l'ordre change}
\end{calculs}
\explique{À la calculatrice (table des valeurs de $P(X\leq k)$) : $P(X\leq 177)\approx
0{,}935$ et $P(X\leq 178)\approx 0{,}959$.}{$0{,}935<0{,}95\leq 0{,}959$}
\explique{$P(X\leq k)$ augmente avec $k$ (on ajoute des probabilités positives) : la
condition est fausse pour $k\leq 177$ et vraie pour $k\geq 178$. Le musée doit préparer
au moins \res{$178$ audioguides}.}{on conclut par une phrase}
\end{correction}

\textbf{Avec Python} (programme de seuil). La fonction \texttt{comb(n, k)} du module
\texttt{math} calcule $\binom{n}{k}$.

\begin{minipage}[t]{0.56\linewidth}\vspace{0pt}
\begin{python}
from math import comb

def seuil(n, p, alpha):
    k = 0
    P = (1 - p)**n
    while P < 1 - alpha:
        k = k + 1
        q = comb(n, k) * p**k * (1 - p)**(n - k)
        P = P + q
    return k
\end{python}
\end{minipage}\hfill
\begin{minipage}[t]{0.4\linewidth}\vspace{0pt}\raggedright
\texttt{P} vaut d'abord $P(X=0)=(1-p)^{n}$, puis on lui ajoute \texttt{q}, qui vaut $P(X=1)$, $P(X=2)$\dots :
à chaque tour, \texttt{P} vaut $P(X\leq k)$.

\smallskip
La boucle s'arrête dès que $P(X\leq k)\geq 1-\alpha$ : \texttt{seuil(200, 0.85, 0.05)}
renvoie $178$.
\end{minipage}

\etape{À vous de jouer}
Une sortie d'observation des étoiles compte $150$ inscrits ; chacun vient avec une
probabilité de $0{,}8$, indépendamment des autres. On admet que le nombre $X$ de
présents suit la loi binomiale $\mathcal{B}(150\,;0{,}8)$. Les organisateurs prévoient
$k$ paires de jumelles. Déterminer le plus petit nombre $k$ de paires de jumelles pour
que la probabilité d'en manquer soit inférieure ou égale à $0{,}05$.
\reponse{$P(X\leq k)\geq 0{,}95$ ; $P(X\leq 127)\approx 0{,}941$ et $P(X\leq 128)\approx
0{,}963$ : $k=128$.}
\end{methodebox}

\afaire{exercices 26 à 29}

\partiecours{Intervalle de probabilité donnée}

Une variable binomiale prend ses valeurs \og autour\fg{} de sa valeur la plus probable.
On cherche un intervalle $\intff{a}{b}$ qui contient $X$ avec une forte probabilité, par
exemple $0{,}95$, en laissant au plus $0{,}025$ de chaque côté.

\begin{proprietebox}[Propriété 7 --- Intervalle de probabilité au moins $0{,}95$]
Soit $X$ une variable aléatoire qui suit une loi binomiale. On note :
\begin{itemize}
  \item $a$ le plus petit entier tel que $P(X\leq a)>0{,}025$ ;
  \item $b$ le plus petit entier tel que $P(X\leq b)\geq 0{,}975$.
\end{itemize}
Alors $P(a\leq X\leq b)\geq 0{,}95$.
\end{proprietebox}

\textbf{\color{coulMeth}Démonstration.} Par définition de $a$, on a $P(X\leq a-1)\leq
0{,}025$ ; et $P(X\leq b)\geq 0{,}975$.
\begin{calculs}
P(a\leq X\leq b) &= P(X\leq b)-P(X\leq a-1) &\qquad& \rem{propriété 5}\\
P(a\leq X\leq b) &\geq 0{,}975-0{,}025 && \rem{on retire au plus $0{,}025$}\\
P(a\leq X\leq b) &\geq 0{,}95 &&
\end{calculs}

\begin{methodebox}[Méthode 9 --- Déterminer un intervalle $I$ tel que $P(X\in I)\geq 0{,}95$]
\etapeexemples
Une enseigne affirme que $80\,\%$ de ses clients sont satisfaits. On interroge au hasard
$100$ clients, en assimilant ce choix à un tirage avec remise, et on note $X$ le nombre
de clients satisfaits. On suppose l'affirmation exacte : $X$ suit la loi binomiale
$\mathcal{B}(100\,;0{,}8)$.
\begin{questions}[label=\alph*)]
  \item Déterminer deux entiers $a$ et $b$ tels que $P(a\leq X\leq b)\geq 0{,}95$, en
        suivant la propriété 7.
  \item Interpréter le résultat.
  \item Sur $100$ clients interrogés, $70$ se disent satisfaits. Qu'en penser ?
\end{questions}
\begin{correction}
\cas{a)}
\explique{À la calculatrice (table des valeurs de $P(X\leq k)$) : $P(X\leq 71)\approx
0{,}020$ et $P(X\leq 72)\approx 0{,}034$, donc \res{$a=72$}.}{plus petit entier tel que
$P(X\leq a)>0{,}025$}
\explique{$P(X\leq 87)\approx 0{,}9747$ et $P(X\leq 88)\approx 0{,}9874$, donc
\res{$b=88$}.}{plus petit entier tel que $P(X\leq b)\geq 0{,}975$}
\explique{D'après la propriété 7, $P(72\leq X\leq 88)\geq 0{,}95$ ; à la calculatrice,
$P(72\leq X\leq 88)\approx 0{,}967$.}{on vérifie}
\cas{b)}
\explique{Si l'affirmation est exacte, il y a au moins $95\,\%$ de chances que le nombre
de clients satisfaits, sur $100$ interrogés, soit \res{compris entre $72$ et
$88$}.}{interprétation dans le contexte}
\cas{c)}
\explique{$70$ n'appartient pas à $\intff{72}{88}$ : une telle observation avait moins de
$5\,\%$ de chances de se produire si l'affirmation était exacte. \res{On peut mettre en
doute l'affirmation} de l'enseigne.}{l'observation sort de l'intervalle}
\end{correction}

\etape{À vous de jouer}
Un musée scientifique affirme que $40\,\%$ de ses visiteurs choisissent l'atelier de
robotique. On interroge au hasard $150$ visiteurs, en assimilant ce choix à un tirage
avec remise, et on note $X$ le nombre de visiteurs qui ont choisi cet atelier. On
suppose l'affirmation exacte : $X$ suit la loi binomiale $\mathcal{B}(150\,;0{,}4)$.
\begin{questions}[label=\alph*)]
  \item Déterminer deux entiers $a$ et $b$ tels que $P(a\leq X\leq b)\geq 0{,}95$, en
        suivant la propriété 7.
  \item Interpréter le résultat.
  \item Sur $150$ visiteurs interrogés, $55$ ont choisi l'atelier de robotique. Qu'en
        penser ?
\end{questions}
\reponse{a) $a=48$ ($P(X\leq 47)\approx 0{,}018$, $P(X\leq 48)\approx 0{,}026$) et $b=72$
($P(X\leq 71)\approx 0{,}972$, $P(X\leq 72)\approx 0{,}981$).\; c)~$55\in\intff{48}{72}$ :
pas de raison de mettre en doute l'affirmation.}
\end{methodebox}

\afaire{exercices 30 à 32, puis les exercices type bac 33 à 35}

\end{document}
